% Generated draft: requires mathematical and source review.
\documentclass[11pt]{article}
\usepackage[margin=25mm]{geometry}
\usepackage{fontspec}
\setmainfont{Latin Modern Roman}
\usepackage{amsmath,amssymb,mathtools,unicode-math}
\setmathfont{Latin Modern Math}
\usepackage{enumitem,needspace,xurl,hyperref}
\setlength{\parindent}{0pt}
\setlength{\parskip}{5pt}
\setlength{\emergencystretch}{3em}
\newcommand{\N}{\mathbb{N}}
\newcommand{\Z}{\mathbb{Z}}
\newcommand{\Q}{\mathbb{Q}}
\newcommand{\R}{\mathbb{R}}
\newcommand{\C}{\mathbb{C}}
\newcommand{\F}{\mathbb{F}}
\newcommand{\A}{\mathbb{A}}
\newcommand{\PP}{\mathbb{P}}
\newcommand{\E}{\mathbb{E}}
\newcommand{\eps}{\varepsilon}
\newcommand{\dd}{\,\mathrm{d}}
\begin{document}
\section*{UnsolvedMath ID: 20002375}
\textbf{Source problem code:} catalogue entry
\subsection{Definitions and mathematical statement}
This source is Elizabeth Brown's participant-interest statement, item A.3 of the appendix ``Participant Contributions,'' not a mathematical conjecture. Brown describes core model theory as largely new to her, with preparation for the workshop, and expresses interest in its relation to large cardinals, their equivalencies, and implications for analysis. No particular cardinal hypothesis, analytic conclusion, or open question is posed. The following is explicitly an auxiliary synthesis supplied by the research aid, not a conjecture attributed to Brown.

Work in ZFC and code reals as elements of \(2^\omega\). For a real \(a\), let \(L[a]\) be the constructible universe relative to \(a\), and let \(\omega_1^{L[a]}\) be its first uncountable ordinal, compared with the ambient \(\omega_1\). A set \(A\subseteq2^\omega\) is \(\Sigma^1_2(a)\) if it has a definition \(x\in A\Longleftrightarrow\exists y\,\forall z\,\theta(x,y,z,a)\) with \(y,z\) real and \(\theta\) arithmetic. It has the perfect set property if it is at most countable or contains a nonempty closed subset with no isolated points.

The aid's parameter-specific bridge is
\[
 \forall a\in2^\omega\quad
 \omega_1^{L[a]}<\omega_1
 \ \Longrightarrow\
 \forall A\subseteq2^\omega\
 \bigl(A\in\Sigma^1_2(a)\Longrightarrow A\text{ has the perfect set property}\bigr).
\]
The supplied argument invokes the relativized Mansfield--Solovay alternative: a \(\Sigma^1_2(a)\) set is contained in \(L[a]\) or contains a perfect subset. Under the displayed cardinal gap, the reals of \(L[a]\) are countable in the ambient universe: internally \(L[a]\) satisfies CH, and its \(\omega_1^{L[a]}\) is externally countable. Thus the first branch of the alternative gives countability.

The aid treats the existence of the canonical sharp \(a^\#\) for \(L[a]\) as a sufficient certificate for \(\omega_1^{L[a]}<\omega_1\), not as an assumption equivalent to the gap. The detailed mouse or indiscernibles convention for sharps is not included in the record. Its all-parameter conclusion is
\[
 \bigl(\forall a\in2^\omega\ a^\#\text{ exists}\bigr)
 \ \Longrightarrow\
 \text{every boldface }\Sigma^1_2\text{ set has the perfect set property}.
\]
Boldface here allows an arbitrary real parameter, so existence of a sharp for just one chosen real supplies only the corresponding parameter-specific implication, not the displayed boldface conclusion. The cardinal-gap premise, the sharp sufficient condition, and the all-real quantifier are separate parts of the supplied synthesis. None changes the context-only nature of the original participant statement.
\subsection{Short English statement}
Brown's source entry states research interests rather than a problem. The aid adds a parameter-specific bridge from \(\omega_1^{L[a]}<\omega_1\) to the perfect set property for \(\Sigma^1_2(a)\) sets, with sharps as a sufficient hypothesis and all real parameters required for the boldface conclusion.
\subsection{Sources}
\begin{itemize}
\item Local registry, permanent UnsolvedMath ID \(20002375\), exact public code \texttt{catalogue entry}; canonical location \texttt{aim-logic-notes.json}, \texttt{notes[150]}.
\item AIM, \emph{Recent advances in core model theory}, appendix ``Participant Contributions,'' A.3, ``Brown, Elizabeth''; supplied reference \url{https://aimath.org/WWN/coremodel/coremodel.pdf}. The verification field identifies A.1--A.4 as participant statements. No source was inspected here. No year, original remarks, or literature entry is supplied.
\item Mansfield--Solovay alternative, invoked by the supplied aid in its relativized parameter form; no full bibliographic citation or theorem number is supplied. The sharp-to-cardinal-gap implication is likewise supplied without a primary citation.
\item Supplied research attempt 1: \texttt{context\_only}, run status \texttt{not\_a\_problem}, rigor label \texttt{proved}, low novelty confidence; registry status \texttt{open}, classification \texttt{OPEN-TRIAGE}, created and updated 14 August 2026. The aid explicitly says it is machine generated, not peer reviewed, and requires independent verification before citation.
\end{itemize}
\end{document}
